\chapter{Related Work}

\section{Review-based recommendation}
Review-based recommender systems use review text to improve rating prediction, preference
modelling or item discovery. Matrix factorisation is a classic collaborative-filtering
baseline, but it mainly models user--item interactions and does not directly explain why an
item satisfies a new textual request \citep{koren2009matrix}. Review-based retrieval addresses
a different problem: the user describes a preference in natural language and the system
retrieves reviewed items whose evidence matches that preference.

The RIRD work is the closest foundation for this project. It defines reviewed-item retrieval
with restaurant reviews, free-text queries and human relevance judgements
\citep{rird2023}. The authors also provide a supervised fusion-based reference, which gives us
a meaningful external comparison. We do not reproduce their fine-tuning; instead, we test
frozen, off-the-shelf retrieval models so that the effect of representation, granularity and
aggregation can be isolated.

\section{Conversational and preference-based recommendation}
Conversational recommender surveys describe systems that elicit preferences through dialogue
and use those preferences to narrow a recommendation set \citep{gao2021conversational}. Work on
preference elicitation also shows the value of asking or generating usage-related questions
when the user's preference is underspecified \citep{kostric2024usage}. Our system is a
single-query, retrieval-focused version of this idea: instead of conducting a conversation, it
accepts one natural-language query and expands it only when the query is vague or contains an
explicit exclusion. This keeps the experiment reproducible and makes the retrieved review
sentence the primary evidence.

\section{Lexical and dense retrieval}
BM25 is the lexical baseline in the RIRD work and in our experiments. It uses term frequency,
inverse document frequency and document-length normalisation
\citep{robertson2009bm25}. BM25 is inexpensive, deterministic and interpretable, but it is
limited by vocabulary mismatch: ``affordable'' and ``cheap'' may describe the same preference
without sharing tokens.

Dense retrieval addresses this limitation by representing queries and text in a shared vector
space. BERT showed how bidirectional Transformer pretraining can provide contextual language
representations \citep{devlin2019bert}; Sentence-BERT then made sentence-level similarity
practical through Siamese encoders \citep{reimers2019sbert}. Dense passage retrieval further
demonstrated the value of independently encoded queries and passages for semantic retrieval
\citep{karpukhin2020dpr}. We use a multilingual sentence-transformer because the same retrieval
procedure must operate on English RIRD text and Thai Wongnai text.

\section{Thai language processing}
Thai differs from English because word boundaries are not consistently marked by spaces.
WangchanBERTa demonstrates the value of Thai-specific pretraining
\citep{lowphansirikul2021wangchanberta}, while our project deliberately uses a single
multilingual encoder to keep the English and Thai experiments comparable. Thai sentence
segmentation and tokenisation are handled with a Thai NLP toolkit.
This is a transfer study rather than a claim that the English RIRD model has been fine-tuned
for Thai.

\section{Granularity, fusion and explanations}
A whole restaurant document preserves context but hides which review supports the result.
Sentence units produce more precise explanations but also create more noise. Our fusion step
therefore compares max, mean and mean-of-top-$k$, with $k\in\{3,5,10\}$. This follows the
reviewed-item retrieval formulation, where unit scores must be converted into an item ranking
\citep{rird2023}.

Explainable recommendation research distinguishes between showing evidence and merely producing
a recommendation that appears plausible. The explainable recommendation survey highlights the
need to connect explanations to user preferences and recommendation decisions
\citep{zhang2020explainable}. Our system selects existing review sentences rather than
generating new claims, and evaluates whether deleting shown sentences affects the ranking.
This makes interpretability a measurable property rather than only a user-interface feature.

\section{Positioning and design justification}
The proposal sources establish the components used here: BM25 for a strong lexical baseline,
Transformer sentence representations for semantic retrieval, and RIRD for reviewed-item
evaluation. We combine them in a deliberately simple pipeline rather than adding a learned
fusion model or a large conversational architecture. The choice is justified by three needs:
fair comparison on the same units, evidence that can be inspected by a diner, and a model that
can be transferred to Thai text without training a separate task-specific system.
